\section{Base cases and symmetric side information}\label{sec:undirected}

This section treats the base case of $2$-dicolourable relations and two classes of graphs: graphs of circular chromatic
number below three, at any degree, and triangle-free graphs of maximum degree three.

\begin{proposition}[$2$-dicolourable relations]\label{prop:twoAcyclic}
\formal{Corollaries.lean}{IsAcyclicIn, rateBound_of_twoAcyclic}
If $V$ is the disjoint union of two sets $A$ and $B$ that are acyclic for $N$, then $N$ satisfies \eqref{eq:star}.
\end{proposition}

\begin{outline}
A global elimination certificate with records $A$ and $B$, each the supplying record of the other: while $A$ is not
revealed, eliminate a vertex of the unrevealed part of $A$ with no parent there; then do the same in $B$. Apply
\cref{prop:global}.
\end{outline}

Bipartite graphs are examples, and so is every loopless relation whose underlying graph is a forest: its directed cycles are
digons, so a proper $2$-colouring of the underlying graph is a $2$-dicolouring. The bound is classical: for a valid $h$,
each of $A$ and $B$ determines the other in topological order, so $h(V)=h(A)=h(B)$.

\paragraph{Why it matters.} This is the base case of every decomposition below: pieces end, after their cyclic parts have
been removed, in $2$-dicolourings. For~(T) it says that side information without odd dependency structure, in this
directed sense, gives no coding gain at all.

\subsection{Circular chromatic number below three}\label{sec:circular}

The \emph{circular clique} $K_{p/q}$ has vertices $\Z/p$, with $i\sim j$ if and only if $q\le(j-i)\bmod p\le p-q$.
The \emph{Andrásfai graph} $\Andr(k)$, $k\ge2$, is $K_{(3k-1)/k}$. The circular chromatic number $\chi_c(G)$ of a finite
graph is the least $p/q$ such that $G$ maps homomorphically to $K_{p/q}$~\cite{Vince1988,Zhu2001}; so $\chi_c(G)<3$ if
and only if $G$ maps to some $K_{p/q}$ with $p<3q$.

\begin{theorem}[Circular chromatic number below three]\label{thm:circular}
\formal{Corollaries.lean}{andrasfaiCert, rateBound_of_circularHom'}
Let $G$ be a finite simple graph with a homomorphism to a circular clique $K_{p/q}$ with $q>0$ and $p<3q$, that is,
with $\chi_c(G)<3$. Then $N_G$ satisfies \eqref{eq:star}.
\end{theorem}

This is the circular memo's theorem~\cite[Corollary~3.2]{CircularMemo2026}, Lean-verified there through the fourth
paper's Theorem~H. We re-derive it as a direct instance of global certificates.

\begin{outline}
The Andrásfai graph $\Andr(k)$ has a global elimination certificate with three records
(\lean{andrasfaiCert}): the intervals $[0,k)$, $[k,2k)$ and $[2k,3k-1)$, each supplied by the next one cyclically.
The elimination order is the circular memo's interleaved order, with keys $3i$, $3j+1$ and $3\ell+2$ for the $i$-th,
$j$-th and $\ell$-th vertex of the three intervals, and the memo's adjacency lemma (\lean{CircularMemo.neighbor_order})
says that every later neighbour of a vertex lies in the next interval. For $k=q+1$ we have $pk\le(3k-1)q$, so the
Bondy--Hell map $i\mapsto\lfloor(3k-1)i/p\rfloor$ is a homomorphism $K_{p/q}\to\Andr(k)$~\cite{BondyHell1990}, formalized
in the fourth paper (\lean{MotionDepthCodes.circularCliqueHom}). Compose it with the given homomorphism, pull the
certificate back (\cref{lem:pullback}) and apply \cref{prop:global}.
\end{outline}

\paragraph{Why it matters.} The theorem has no degree assumption: every graph of circular chromatic number below three,
for instance every odd cycle of length at least five and every graph mapping to one, satisfies the half bound for all
codes over all alphabets. For~(T), symmetric side information of this kind cannot carry a coding gain however dense it
is. The bound does not reach $\chi_c=3$; the Petersen graph, which has $\chi_c=3$~\cite{JHS1997}, is covered by the
fourth paper's certificate~\cite[Corollary~6.9]{FourthPaper2026} and by the next theorem.

\subsection{Triangle-free graphs of maximum degree three}\label{sec:N3}

\begin{theorem}[Theorem N3]\label{thm:N3}
\formal{N3.lean}{CycIn, CycIn.chordless, CycIn.share, cyclePiece, decomp_of_cycles, n3_decomp, rateBound_n3, card_sq_le_n3}
Let $G$ be a triangle-free graph of maximum degree at most~$3$. Then $N_G$ has a decomposition. Hence $N_G$ satisfies
\eqref{eq:star}, and every storage code $C$ for $G$ over a nonempty finite alphabet satisfies
$\card C^2\le\card\alpha^{\card V}$.
\end{theorem}

\begin{outline}
By strong induction on the set $R$ of remaining vertices, with context $V\setminus R$, \lean{decomp_of_cycles}
constructs a decomposition from every context. If $G[R]$ contains a cycle, let $C=c_0c_1\cdots c_{g-1}$ be a shortest one
(\lean{CycIn}); $g\ge4$ because $G$ is triangle-free. Three facts:
\begin{enumerate}[label=(\roman*)]
\item $C$ is chordless (\lean{CycIn.chordless}), since a chord would close a shorter cycle;
\item if a vertex $w\in R$ off $C$ is adjacent to $c_i$ and $c_j$, $i<j$, it closes cycles of lengths $j-i+2$ and
$g-(j-i)+2$, both at least $g$; hence $g=4$ and $c_i,c_j$ are opposite (\lean{CycIn.share});
\item every $c_i$ has at most one neighbour off $C$, since it has two on $C$ and degree at most three.
\end{enumerate}
The piece certificate of $C$ in context $V\setminus R$ (\lean{cyclePiece}) eliminates $c_0,\dots,c_{g-1}$ in order with
three records given by colours: $c_0$ has colour~$0$ and $c_{g-1}$ colour~$1$; if $g$ is odd, $c_{g-2}$ has
colour~$2$; every other $c_i$ has colour~$1$ for odd $i$ and~$0$ for even $i$. The supplying record of $c_t$ is the
colour of $c_{t+1}$, and that of $c_{g-1}$ is colour~$0$. At a step $t\ge1$ the cycle neighbours of $c_t$ are
$c_{t-1}$, already eliminated, and $c_{t+1}$, in the supplying record; at step~$0$ they are $c_1$ and $c_{g-1}$, both of
colour~$1$. The remaining neighbour of $c_t$, if it lies in $R\setminus C$, is its outside parent. By~(ii) outside
parents are private, except for opposite vertices of a $4$-cycle, which have equal supplying records ($c_0,c_2$ are
supplied by colour~$1$, and $c_1,c_3$ by colour~$0$); \cref{lem:cosupply} applies. If $G[R]$ is a forest, a proper
$2$-colouring of it is a piece by \cref{lem:Bdir}. Theorem~D concludes.
\end{outline}

By \cref{lem:fewer-arcs}, the half bound holds also for every relation $N$ with $N(v)\subseteq N_G(v)$ for all $v$,
that is, for every relation whose arcs join adjacent vertices of $G$, in one direction or in both;
\lean{decomp_of_cycles} gives decompositions for such relations as well, its hypothesis on $3$-cycles being vacuous here.
Several known results are special cases.
\begin{itemize}[leftmargin=2em]
\item The fourth paper's Theorem~S~\cite[Theorem~6.13]{FourthPaper2026} proves the half bound for scalar linear codes,
over every field, on a directed class that contains these graphs; Theorem~N3 is its all-code version on triangle-free
subcubic graphs. For scalar linear codes the exact value is $\card V-\nu(G)$ for the minrank, with the integral matching
number $\nu$~\cite[Theorem~6.16]{FourthPaper2026}; for all codes the guessing number is the fractional matching number
$\nustar(G)$ (\cref{thm:gn}). The two differ already on $C_5$, where $\nu=2$ and $\nustar=5/2$.
\item The subcubic memo~\cite[Theorem~7.1]{SubcubicMemo2026} proves the half bound, in Lean, for triangle-free subcubic
graphs all of whose edges lie in four-cycles, through a homomorphism to $K_{8/3}$; this is a special case. Its
information-set theorem~\cite[Theorem~2.3]{SubcubicMemo2026} bounds codes that are reversible on the rows by
$q^{\nu(G)}$, which is stronger on that class of codes.
\item The circular memo's theorem (\cref{thm:circular}) covers graphs of every degree with $\chi_c<3$, and the fourth
paper's certificates cover the graphs that map to the Petersen graph~\cite[Corollary~6.9]{FourthPaper2026}. Theorem~N3
needs no target: it covers every triangle-free subcubic graph. In particular it proves unconditionally the fourth paper's
conditional result for cubic graphs of large girth~\cite[Corollary~6.10]{FourthPaper2026}, which assumed that such graphs
map to $C_5$ or to $K_{14/5}$.
\item Dang~\cite{Dang2015} computes the guessing numbers of several families of planar triangle-free graphs and
conjectures that the fractional clique cover strategy is optimal on every planar triangle-free
graph~\cite[Conjecture~15]{Dang2015}. With \cref{thm:gn} below, Theorem~N3 proves this for maximum degree three,
without planarity.
\end{itemize}

\paragraph{Why it matters.} No storage code on a triangle-free subcubic graph beats the pairing, whatever its alphabet
and whether or not it is linear. With \cref{cor:paid}, this answers the fourth paper's undirected
crux~\cite[Question~6.1]{FourthPaper2026} affirmatively at maximum degree three: in a locally bipartite family of maximum
degree at most three, only $o(\card V)$ vertices lie on triangles, deleting them leaves a triangle-free subcubic graph, and
so every storage code has rate at most $\frac12+o(1)$. For~(T), symmetric subcubic side
information that is triangle-free cannot carry a coding gain. Native cone relations are directed, however, and their
symmetrizations can contain linearly many triangles~\cite[Proposition~3.10]{DirectedRecovery2026}; this is why
\cref{sec:directed} allows directed triangles.
